\part{Language Support \& Development Tooling}
\label{part:lang}

With packages loaded and configured (Part~\ref{part:packages}), the
question becomes: what do they actually \emph{do} when you open a
source file?  This part traces the path from file extension to
fully-powered editing session.

Chapter~\ref{ch:tree-sitter} opens with the two mode-dispatch
mechanisms---\texttt{auto-mode-alist} and
\texttt{major-mode-remap-alist}---that select between legacy and
tree-sitter grammars.
Chapter~\ref{ch:major-minor-modes} steps back to explain the
major/minor-mode infrastructure that those dispatchers target.
Chapters~\ref{ch:eglot-lsp} and~\ref{ch:eglot-server-programs}
cover the LSP integration layer: how Eglot maps protocol capabilities
to Emacs interfaces, and how to register custom server contacts
without falling into the \texttt{mapcan} circular-list trap.
Chapter~\ref{ch:find-file-visit-truename} addresses a subtle
symlink-resolution issue that can break LSP buffer identity.
Chapter~\ref{ch:dape} adds the runtime counterpart to LSP: the
Debug Adapter Protocol.  Chapter~\ref{ch:fringe-bitmaps} explains
the fringe display model that git-gutter to fold-indicators rely on.
The part closes with Chapter~\ref{ch:testing}, which explains how
ERT and Buttercup guard the configuration against regressions---
ensuring that performance changes from Part~\ref{part:startup} do
not silently break functionality.

\chapter{Tree-sitter mode selection}
\label{ch:tree-sitter}
\minitoc\label{tree-sitter-mode-selection}
\index{tree-sitter}\index{auto-mode-alist}\index{major-mode-remap-alist}\index{treesit-auto}

% ── Opening signpost ──────────────────────────────────────────────
Tree-sitter provides incremental parsing for syntax highlighting,
folding, and structural navigation, but activating it requires
understanding the difference between \texttt{auto-mode-alist} and
\texttt{major-mode-remap-alist}.  This chapter explains both dispatch
mechanisms and why the selective-remap pattern is vastly more
efficient than \texttt{global-treesit-auto-mode}.  The mode
infrastructure that these dispatchers target is covered in
Chapter~\ref{ch:major-minor-modes}.

\section{\texttt{auto-mode-alist} vs \texttt{major-mode-remap-alist}}\label{auto-mode-alist-vs-major-mode-remap-alist}

These are two separate mechanisms and understanding the distinction is key:

\begin{itemize}
\item \textbf{\texttt{auto-mode-alist}}: Maps filename patterns to major modes. \texttt{".py"} → \texttt{python-mode}. This is the \textbf{original} dispatch.
\item \textbf{\texttt{major-mode-remap-alist}} (Emacs 29+): After \texttt{auto-mode-alist} selects a mode, this alist can \textbf{remap} it to a different mode. \texttt{python-mode} → \texttt{python-ts-mode}. Zero per-file overhead: it's a simple alist lookup.
\end{itemize}

\section{Why \texttt{global-treesit-auto-mode} is expensive}\label{why-global-treesit-auto-mode-is-expensive}

\texttt{treesit-auto} provides \texttt{global-treesit-auto-mode}, which automatically remaps modes to their tree-sitter variants. Under the hood, it calls \texttt{treesit-language-available-p} for \textbf{every} language it knows about (\(\sim\)30+) to build the remap list. Each \texttt{treesit-language-available-p} call probes the filesystem for a shared library (\texttt{libtree-sitter-*.so}).

On macOS with \(\sim\)30 languages, this costs \(\sim\)1.8 s per file-open (the probe runs on \texttt{after-change-major-mode-hook}, which fires every time a major mode activates).

\section{The selective-remap pattern (what your config uses)}\label{the-selective-remap-pattern-what-your-config-uses}

Instead of the expensive global mode, explicitly remap only the languages you use:

\begin{minted}{elisp}
(dolist (pair '((python-mode     . python-ts-mode)
                (java-mode       . java-ts-mode)
                (js-json-mode    . json-ts-mode)
                (css-mode        . css-ts-mode)
                (html-mode       . html-ts-mode)
                (mhtml-mode      . html-ts-mode)
                (dockerfile-mode . dockerfile-ts-mode)))
  (add-to-list 'major-mode-remap-alist pair))
\end{minted}

For modes that have \textbf{no} non-ts entry in the default \texttt{auto-mode-alist} (e.g., YAML, Bash), add direct filename associations:

\begin{minted}{elisp}
(add-to-list 'auto-mode-alist '("\\.ya?ml\\'" . yaml-ts-mode))
(add-to-list 'auto-mode-alist '("\\.sh\\'"    . bash-ts-mode))
\end{minted}

Keep \texttt{treesit-auto} around with \texttt{:demand t} \textbf{only} for its grammar installation command (\texttt{treesit-auto-install}), not for its global mode.

% ── Mental model summary ──────────────────────────────────────────
\section{Mental model}
Emacs resolves a buffer's major mode in two stages:
\texttt{auto-mode-alist} maps filename patterns to modes, and
\texttt{major-mode-remap-alist} can then redirect the chosen mode to an
alternative (typically its tree-sitter variant) after the initial
dispatch.  Enabling \texttt{global-treesit-auto-mode} short-circuits
this by probing the filesystem for \(\sim\)30 shared grammar libraries
on every mode activation, which is expensive.  The selective-remap
pattern avoids that cost entirely: you register only the languages you
actually use, so each file open incurs zero additional overhead.  For
modes that lack a non-tree-sitter \texttt{auto-mode-alist} entry (YAML
and Bash are common examples), you can add direct filename associations
instead of relying on remapping.  The recommended approach is to keep
\texttt{treesit-auto} installed solely for its grammar installation
command and leave its global mode disabled.
